\section{Linear slices and balanced endpoint transfer}\label{sec:slices}
Fix $A_1,\ldots,A_{m-1}\in\Mat_n(K)$ and set
\begin{equation}\label{eq:slice}
 L_A(X)=f(A_1,\ldots,A_{m-1},X).
\end{equation}
For multilinear $f$, the map $L_A$ is linear. Thus any linear
combination of outputs obtained by varying $X$ \emph{with the same fixed
tuple $A$} is again one value of $f$. No additive closure of the
unrestricted image is assumed. Simultaneous conjugation of all inputs
conjugates the output of $s_m$.

\begin{lemma}[Jordan-path criterion]\label{lem:jordanpath}
Suppose a slice image contains all diagonal matrices and the matrix
units
\[E_{i_1i_2},E_{i_2i_3},\ldots,E_{i_{n-1}i_n}\]
for some permutation $(i_1,\ldots,i_n)$ of the indices. Then
$\im f(\Mat_n(K))=\Mat_n(K)$, provided $f$ is equivariant under
simultaneous conjugation.
\end{lemma}
\begin{proof}
Within this one slice we can prescribe each diagonal entry and the
superdiagonal entries along the indicated path independently. A
Jordan matrix, including a direct sum of several blocks, has precisely
this form after permuting the coordinates, with its link entries
chosen to be $1$ or $0$. Every matrix over $K$ is similar to such a
matrix. Conjugating the entire input tuple concludes the proof.
\end{proof}

For later use let
\begin{equation}\label{eq:V2}
 V_h=\Span\{E_{i,i+h\bmod n}:0\le i<n\}.
\end{equation}
If $n$ is odd, the step-two edges in $V_2$ contain a Hamiltonian
cycle. If $n$ is even, the two parity cycles are linked by $E_{03}$;
an explicit Hamiltonian path is
\begin{equation}\label{eq:evenpath}
 2\longrightarrow4\longrightarrow\cdots\longrightarrow n-2
 \longrightarrow0\longrightarrow3\longrightarrow5
 \longrightarrow\cdots\longrightarrow n-1\longrightarrow1.
\end{equation}

Let $f$ be homogeneous multilinear of degree $m$ and let its final
variable be free. In the signed word expansion, denote by $P(A)$
the total fixed-input product coefficient when that variable is at
the \emph{last} position; let $Q(A)$ be its counterpart in the
\emph{first} position. Both depend polynomially on $A$.

\begin{theorem}[Balanced endpoint transfer]\label{thm:transfer}
Let $n=p+q$ with $|p-q|\le1$. Suppose there are fixed
$(m-1)$-tuples $A\in\Mat_p(K)^{m-1}$ and
$B\in\Mat_q(K)^{m-1}$ such that both slice images contain every
diagonal matrix in their respective blocks and $P(A)$ and $Q(A)$
are invertible. Then $\im f(\Mat_n(K))=\Mat_n(K)$, assuming
simultaneous-conjugacy equivariance.
\end{theorem}
\begin{proof}
Take the block-diagonal fixed tuple
\[
 C_j(t)=\begin{pmatrix}A_j&0\\0&tB_j\end{pmatrix}
 \qquad(1\le j<m).
\]
Writing the free input in four rectangular blocks yields four
independent block-preserving linear maps. On the two diagonal blocks
they are $L_A$ and $t^{m-1}L_B$. On a top-right input $X$,
setting $t=0$ kills terms in which a $B$ factor occurs on the
right of $X$, giving $X\mapsto P(A)X$. On a bottom-left input, the
limit is $X\mapsto XQ(A)$. These last maps are invertible for $t=0$.
The corresponding determinants are nonzero polynomials in $t$.
Choose a single $t\in K^\times$ avoiding their finitely many zero
sets. This choice is independent of the eventual target matrix.

At this fixed nonzero $t$, the slice covers all four mixed-block
entries and all diagonal entries in either square block. Since
$|p-q|\le1$, the bipartite graph of entries between the two blocks
has a Hamiltonian path alternating between the blocks. Therefore
Lemma~\ref{lem:jordanpath} applies.
\end{proof}

For $f=s_{2k+1}$ the two endpoint coefficients are
\begin{equation}\label{eq:endpoints}
 P(A)=Q(A)=s_{2k}(A_1,\ldots,A_{2k}).
\end{equation}
Indeed moving the free letter across the other $2k$ letters changes
no sign. The balanced hypothesis is essential: no Hamiltonian path
through a complete bipartite graph exists if its two parts differ in
size by more than one.

\section{Signed trails on a bidirected path}\label{sec:path}
Fix $k\ge1$, write $r=k+1$, and consider the ordered list
\begin{equation}\label{eq:pathlist}
 \mathcal A_k=(E_{01},E_{12},\ldots,E_{k-1,k},
 E_{k,k-1},E_{k-1,k-2},\ldots,E_{10}).
\end{equation}
Set
\[
\begin{aligned}
 c&=(1,\underbrace{2,\ldots,2}_{r-2},1)^t,\\
 H&=\diag(c)+cc^t,\\
 \mathcal E_r&=\diag(1,-1,1,\ldots,(-1)^{r-1}).
\end{aligned}
\]
For $r=2$ the middle list is empty. Every nonzero product of
matrix units is an oriented Euler trail through their directed
edges. This simple observation gives the exact slice coefficients.

\begin{lemma}[Endpoint and loops]\label{lem:pathloop}
For the tuple \eqref{eq:pathlist},
\begin{align}
 s_{2k}(\mathcal A_k)&=\diag((-1)^ic_i)_{0\le i<r},\label{eq:pathendpoint}\\
 s_{2k+1}(\mathcal A_k,\diag(x))&=\diag(H\mathcal E_r x).
 \label{eq:pathdiag}
\end{align}
Furthermore,
\begin{equation}\label{eq:pathdet}
 \det(H\mathcal E_r)=(-1)^{r(r-1)/2}2^{r-2}(2r-1)\ne0.
\end{equation}
\end{lemma}
\begin{proof}
The forward and reverse directed edges form a bidirected path. Every
closed Euler ordering is a cyclic rotation of the tour traversing
that path forward and returning. Vertex $i$ occurs as a starting
vertex $c_i$ times, and the resulting permutations have sign
$(-1)^i$. This proves \eqref{eq:pathendpoint}.

Insert a diagonal loop at vertex $j$. For a tour beginning at $i$,
there are $c_j+\delta_{ij}$ loop slots. Incorporating the loop into
the fixed ordering gives sign $(-1)^j$ for the added-loop
coefficient, so its contribution to output coordinate $i$ is
$c_i(c_j+\delta_{ij})(-1)^j$. This is exactly
$H\mathcal E_r$. Finally the determinant lemma yields
\[
 \det H=\left(\prod_{i=0}^{r-1}c_i\right)
 \left(1+c^t\diag(c)^{-1}c\right)
 =2^{r-2}(2r-1).
\]
Multiply by $\det\mathcal E_r=(-1)^{r(r-1)/2}$.
\end{proof}

\Needspace{18\baselineskip}
\begin{lemma}[Free chord on a bidirected path]\label{lem:pathchord}
Embed \eqref{eq:pathlist} on an interval of indices
$a,a+1,\ldots,a+k$, with $c_j$ as above. For $p\ne q$, the
image of the free matrix unit $E_{pq}$ is a multiple of $E_{pq}$.
The coefficient, before any external window weight, is
\begin{equation}\label{eq:pathchordtable}
\begin{array}{@{}ll@{}}
\text{position of $p,q$}&\text{coefficient}\\\midrule
p,q\text{ inside, same parity}&(-1)^{p-a}c_{p-a}c_{q-a}\\
p,q\text{ inside, opposite parity}&0\\
p\text{ inside},\ q\text{ outside}&(-1)^{p-a}c_{p-a}\\
q\text{ inside},\ p\text{ outside}&(-1)^{q-a}c_{q-a}\\
p,q\text{ outside}&0.
\end{array}
\end{equation}
\end{lemma}
\begin{proof}
The directed degree imbalance of the fixed bidirected path is zero.
Adding $p\to q$ forces every connected Euler trail to start at $p$
and end at $q$, so the output unit is $E_{pq}$. If the free edge
joins a path vertex to an external vertex, it must be first or
last; \eqref{eq:pathendpoint} gives the two one-sided coefficients.
If both endpoints lie outside, the selected edges are disconnected.

For two inside vertices, assume first $p<q$ and put $d=q-p$.
The tails outside $[p,q]$ must be traversed outward and back. Each
nonempty tail admits the independent attachment multiplicity
$c_{p-a}$ or $c_{q-a}$. After stripping these tails, an Euler trail
is specified by a turning vertex $j\in[p,q]$: go from $p$ to $j$
and back, take the chord $p\to q$, then go from $q$ to $j$ and
back. Moving $j$ one edge changes the permutation sign by $-1$.
The initial sign, by inversion counting in the displayed order of
\eqref{eq:pathlist}, is $(-1)^{q-a}$. Summing gives
\[
 (-1)^{q-a}c_{p-a}c_{q-a}
 \sum_{h=0}^{d}(-1)^h.
\]
The sum vanishes for odd $d$ and equals $1$ for even $d$, in
which case $p$ and $q$ have the same parity. For $p>q$, reversing
both the directed path and the fixed ordering contributes the
same sign $(-1)^k$ twice, giving the identical formula.
\end{proof}

\section{Positive path cores and large matrix sizes}\label{sec:large}
We next build one slice with invertible diagonal restriction and
invertible even-standard endpoint on every $\Mat_q(K)$ with $q\ge r$.
For $k=r-1$, define
\begin{equation}\label{eq:Vandermonde}
\begin{split}
 F_j&=\sum_{e=0}^{q-2}(e+1)^{j-1}E_{e,e+1},\\
 R_j&=\sum_{e=0}^{q-2}z_e(e+1)^{j-1}E_{e+1,e},
 \qquad z_e=(-1)^{\lfloor e/k\rfloor},
\end{split}
\end{equation}
for $1\le j\le k$, and order the fixed inputs as
$F_1,\ldots,F_k,R_k,\ldots,R_1$. Put
\[
 v_k=\prod_{j=1}^{k-1}j!,\qquad
 \iota_a:K^r\longrightarrow K^q,
 \quad(x_0,\ldots,x_{r-1})\longmapsto
 \sum_{i=0}^{r-1}x_i e_{a+i}.
\]

\begin{proposition}\label{prop:positive}
For every $q\ge r$, the diagonal restriction of the final-variable
slice associated with \eqref{eq:Vandermonde} is
\begin{equation}\label{eq:Gram}
 L_{\rm diag}=G\mathcal E_q,
 \qquad G=v_k^2\sum_{a=0}^{q-r}\iota_aH\iota_a^t.
\end{equation}
Its fixed-input even standard polynomial is
\begin{equation}\label{eq:largeendpoint}
 s_{2k}(F_1,\ldots,F_k,R_k,\ldots,R_1)
 =v_k^2\diag\left((-1)^i
 \sum_{\substack{0\le a\le q-r\\a\le i\le a+k}}
 c_{i-a}\right)_{0\le i<q}.
\end{equation}
Both matrices are invertible in characteristic zero.
\end{proposition}
\begin{proof}
Expand the fixed inputs into path units. Alternation among the
$F_j$ requires $k$ distinct forward edges, and alternation among
the $R_j$ requires $k$ distinct backward edges. A diagonal output
has zero net flow through every edge cut, so the forward and reverse
edge sets coincide. With only one extra diagonal loop, a disconnected
set of edges cannot support one Euler trail. Consequently the $k$
selected edges form a consecutive interval
$a,a+1,\ldots,a+k-1$.

On a fixed window, the coefficients from the forward and reverse
choices are two Vandermonde determinants, both of absolute algebraic
value $v_k$, before reverse signs. Their product is $v_k^2$.
Furthermore
\[
 \prod_{e=a}^{a+k-1}z_e=(-1)^a:
\]
the product is $+1$ at $a=0$, and shifting the window changes its
sign because $z_{a+k}/z_a=-1$. Multiplying the local formula
\eqref{eq:pathdiag} by this sign converts the local alternating
factor into the restriction of $\mathcal E_q$. Thus the windows
sum to \eqref{eq:Gram}; \eqref{eq:pathendpoint} similarly yields
\eqref{eq:largeendpoint}.

There is no assumption that $K$ is ordered. Instead regard these
integer matrices first over $\mathbb R$. The matrix
$H=\diag(c)+cc^t$ is positive definite, as is each of its
nonzero embedded quadratic forms. Every index $i$ occurs in at least
one interval $[a,a+k]$, so for $x\ne0$,
\[
 x^tGx=v_k^2\sum_{a=0}^{q-r}
 (\iota_a^tx)^tH(\iota_a^tx)>0.
\]
Therefore $\det G$ is a nonzero integer, and remains nonzero over
any field of characteristic zero. All diagonal entries in
\eqref{eq:largeendpoint} are nonzero signed positive integers.
\end{proof}

\begin{corollary}\label{cor:large}
For $r\ge2$ and $n\ge2r$, the image of $s_{2r-1}$ on
$\Mat_n(K)$ is all of $\Mat_n(K)$.
\end{corollary}
\begin{proof}
Let $p=\lceil n/2\rceil$, $q=\lfloor n/2\rfloor$.
Both are at least $r$; apply Proposition~\ref{prop:positive}
to each block and Theorem~\ref{thm:transfer}, using
\eqref{eq:endpoints}.
\end{proof}

\section{Minimal size: the cycle and its maximal-minor bivector}
\label{sec:minimal}
We treat the boundary
\begin{equation}\label{eq:minimalclaim}
 \im s_{2n-1}(\Mat_n(K))=\Mat_n(K),\qquad n\ge3.
\end{equation}
At this size, a single fixed slice need not serve all targets. The
construction instead adjusts the diagonal fixed inputs to each
prescribed Jordan matrix.

Fix the $n$ directed cycle edges
\[
 C_i=E_{i,i+1\bmod n},\qquad 0\le i<n,
\]
followed by $n-2$ diagonal matrices
$D_1,\ldots,D_{n-2}$ with diagonal row matrix
$D\in\Mat_{n-2,n}(K)$. Leave the final variable free and denote
the resulting slice by $L_D$. Let
\begin{equation}\label{eq:An}
 A_n=2\one\one^t-I_n.
\end{equation}
For a row-deleted coordinate matrix of rank $n-2$, the signed
maximal minors of $D$ form a skew matrix $B(D)$, described
explicitly in the following lemma. Its scale is fixed by the
order of the cycle inputs followed by the diagonal inputs.

\begin{lemma}[Cycle--bivector formulas]\label{lem:cycle}
For $0\le i<j<n$, let $D_{\widehat{i}\widehat{j}}$ be the
square submatrix obtained by removing columns $i,j$, and put
\begin{equation}\label{eq:gamma}
\begin{aligned}
 \gamma_{ij}&=(-1)^{n(n+1)/2+j-i-2},\\
 B_{ji}&=\gamma_{ij}\det D_{\widehat{i}\widehat{j}},
 \quad B_{ij}=-B_{ji},\quad B_{ii}=0.
\end{aligned}
\end{equation}
Then
\begin{equation}\label{eq:cycleDiag}
 L_D(\diag(x))=\diag(A_nB x).
\end{equation}
If $D$ has rank $n-2$, then $B$ has rank two and is decomposable.
Conversely, for independent $v,w\in K^n$, choosing the rows of
$D$ to span $\{x:v^tx=w^tx=0\}$ makes
$B$ a nonzero scalar multiple of $vw^t-wv^t$.

For $p\ne q$, let $I(p,q)$ be the inclusive forward cyclic
interval $q,q+1,\ldots,p$ and define
\begin{equation}\label{eq:ab}
 a_i=\begin{cases}
 1,&i=p\text{ or }i=q,\\2,&i\in I(p,q)\setminus\{p,q\},\\0,&i\notin I(p,q),
 \end{cases}\qquad
 b_i=\mathbf1_{i\notin I(p,q)}.
\end{equation}
The free chord $E_{pq}$ obeys the exact formula
\begin{equation}\label{eq:cycleChord}
 L_D(E_{pq})=-2(a^tBb)E_{pq}.
\end{equation}
In particular the original forward cycle edges have zero
coefficient.
\end{lemma}
\begin{proof}
All coefficients are alternating multilinear expressions in the
$n-2$ diagonal rows of $D$. It is therefore enough to calculate them
on coordinate diagonal inputs omitting a pair of vertices $i<j$;
expansion in minors then gives the stated general formulas.

For a free loop at $i$, the cycle and the $n-2$ fixed loops admit
closed Euler orderings obtained by rotating a cyclic ordering.
The vertex $j$ without a loop occurs once as a starting point and
all other vertices twice. Since the total word length is $2n-1$,
these rotations have the same sign. Place the diagonal loops in
increasing coordinate order after the cycle edges. Inversion
counting of the rotation and its loop insertions gives the common
factor
\[
 \gamma_{ij}=(-1)^{n(n+1)/2+j-i-2}.
\]
Thus the free loop at $i$ produces
$\gamma_{ij}\diag(2\one-e_j)$; the loop at $j$ gives the
negative corresponding expression with $i$ omitted. The two
columns are exactly $A_nB_{\cdot i}$ and $A_nB_{\cdot j}$.
This establishes \eqref{eq:cycleDiag} in the coordinate case.

For completeness we record the chord sign calculation in the same
minor normalization. An Euler trail containing the free chord
$p\to q$ has two possible edge skeletons before the fixed loops
are inserted: chord first followed by the forward cycle starting
at $q$, or the cycle from $p$ followed by the chord. A diagonal
loop strictly in the inclusive interval $I(p,q)$ has a forced
placement, while at either chord endpoint there are two placements.
Pairing the skeletons for a pair of omitted loops on the same side
of the interval cancels their signs. Moving an omitted loop
across an endpoint changes the inversion parity; the resulting
complete coefficient table, for the elementary bivector
$B_{ij}=-B_{ji}$, is
\begin{equation}\label{eq:chordsigns}
\begin{array}{@{}lc@{}}
\text{positions of omitted }i,j\ (i<j)&\text{coefficient of }B_{ij}\\\hline
 i,j\in I\text{ or }i,j\notin I&0\\
 i\text{ endpoint of }I,\ j\notin I&-2\\
 i\text{ strictly inside }I,\ j\notin I&-4\\
 i\notin I,\ j\text{ endpoint of }I&2\\
 i\notin I,\ j\text{ strictly inside }I&4.
\end{array}
\end{equation}
For example, take $n=4$, $p=2$, $q=0$ and omit the loop vertices
$i=0,j=3$, leaving $D_1=E_{11}$ and $D_2=E_{22}$.
Writing $X=E_{20}$ and keeping the input order
$C_0,C_1,C_2,C_3,D_1,D_2,X$, the four admissible Euler words and
permutation signs are
\[
\begin{array}{@{}lr@{}}
 C_2C_3C_0D_1C_1D_2X & -1\\
 XC_0D_1C_1D_2C_2C_3 & -1\\
 D_2C_2C_3C_0D_1C_1X & +1\\
 D_2XC_0D_1C_1C_2C_3 & -1.
\end{array}
\]
Their sum is $-2$, as in the endpoint/outside row of the table:
\eqref{eq:gamma} gives $\gamma_{03}=-1$, hence $B_{03}=1$.
Here is a parameter-uniform inversion check of the table, including
cyclic wraps. Write $S=\{0,\ldots,n-1\}\setminus\{i,j\}$, and
let $\ell=(p-q)\bmod n\in\{1,\ldots,n-1\}$. For the first
skeleton, with any endpoint loop placed at its earliest opportunity,
list the vertices in the order
\[
 \pi_{p,q}=(p,q,q+1,\ldots,p-1,p+1,\ldots,q-1),
\]
where each interval runs in the forward cyclic direction. Deleting
$i,j$ gives the loop ordering $\pi_{p,q}|_S$. Define
$\eta(p)=\eta(q)=n$ and
$\eta(v)=n-((v-q)\bmod n)$ otherwise, and write
$\operatorname{inv}$ for the number of numerical inversions.
The sign of this baseline Euler word is
\begin{equation}\label{eq:cycle-uniform-sign}
 \varepsilon=(-1)^{q(n-q)+\mathbf1_{p\in S}
       +\operatorname{inv}(\pi_{p,q}|_S)
       +\sum_{v\in S}\eta(v)}.
\end{equation}
The terms count, respectively, the rotation of the labelled cycle,
the last-labelled chord crossing the optional early $p$-loop,
inversions among the fixed-loop labels, and the cycle letters
following their respective loops. In particular, the endpoint
values of $\eta$ cannot be replaced by its interior formula.

To see that this absolute sign matches the cofactor convention, let
$z_v$ be the zero-based position of $v$ in $\pi_{p,q}$.
Moving $p$ from position $\ell$ to the front of the forward
rotation starting at $q$ yields
\[
 \operatorname{inv}(\pi_{p,q})\equiv q(n-q)+\ell,
 \qquad \sum_{v=0}^{n-1}\eta(v)=\frac{n(n+1)}2+\ell
 \pmod2.
\]
Deleting $v$ changes inversion parity by $z_v+v$; deleting both
$i<j$ requires the additional parity correction
$\mathbf1_{z_j<z_i}$. Substitution in
\eqref{eq:cycle-uniform-sign} gives
\[
 \frac{\varepsilon}{\gamma_{ij}}=
 (-1)^{1+h(i)+h(j)+\mathbf1_{z_j<z_i}},
 \quad h(v)=z_v+\eta(v)+\mathbf1_{v=p}\pmod2.
\]
The vertex order gives $h(v)\equiv n+1$ for $v\in I(p,q)$ and
$h(v)\equiv n$ outside. When $i,j$ lie on opposite sides,
$h(i)+h(j)\equiv1$, and all interval vertices precede the
outside vertices in $\pi_{p,q}$. Hence
\begin{equation}\label{eq:cycle-uniform-normalization}
 \varepsilon=\gamma_{ij}\operatorname{sgn}(a_i b_j-a_j b_i)
 \quad\text{whenever }a_i b_j-a_j b_i\ne0.
\end{equation}

For clarity, the remaining sign summation is finite independently
of $n$. If $A$ is the forward arc from $q$ to $p$ and $B$ its
complement, their edge lengths, including forced interior loops,
are $\alpha=2\ell-1-u$ and $\beta=2(n-\ell)-1-v$,
where $u,v$ count the omitted strict-interior loops in the
respective arcs. Let $P,Q\in\{0,1\}$ indicate the presence of
endpoint loops at $p,q$. In the first skeleton the parity changes
on moving a present $p$-loop or $q$-loop to its late placement
are $1+Q+\alpha$ and $\alpha+\beta$; in the second they are
$\alpha+\beta+Q$ and $\alpha+1$. The baseline skeletons differ
by $\beta(1+Q+\alpha)+Q+\alpha$ inversions.
Summing the at most $2^{P+Q+1}\le8$ possible words therefore
makes both same-side configurations cancel. In an opposite-side
configuration the signed sum divided by $\varepsilon$ is $2$
when the omitted interval vertex is an endpoint, and $4$
when it is a strict interior vertex. Together with
\eqref{eq:cycle-uniform-normalization}, the complete sum is
$2\gamma_{ij}(a_i b_j-a_j b_i)$, namely the table
\eqref{eq:chordsigns} with $B_{ij}=-\gamma_{ij}$.
Expansion in the maximal minors proves \eqref{eq:cycleChord}.
The degree imbalance of the chord forces output $E_{pq}$,
so no other matrix units occur.

The maximal-minor cofactor construction is the Pluecker bivector
of the two-dimensional annihilator of the row space of $D$.
For rank $n-2$, it is nonzero and decomposable, and its image is
that annihilator. In particular taking the row space to be
$\ker v^t\cap\ker w^t$ gives the claimed multiple of
$vw^t-wv^t$. No positivity or nonisotropy assumption on $v,w$
is involved.
\end{proof}

\begin{proposition}[Every Jordan target at the minimal size]
\label{prop:minJordan}
For each $n\ge3$ and every lower Jordan matrix $J\in\Mat_n(K)$,
there is a rank-$n-2$ choice of $D$ such that $J$ is contained in
the image of the single slice $L_D$, after a possible simultaneous
permutation of coordinates.
\end{proposition}
\begin{proof}
Write $d\in K^n$ for the diagonal of $J$ and $\tau=\sum_i d_i$.
The matrix $A_n$ is invertible: it acts by $2n-1$ on the line
$K\one$ and by $-1$ on its sum-zero complement. If $d\ne0$, put
\begin{equation}\label{eq:vdiag}
 v=A_n^{-1}d=-d+\frac{2\tau}{2n-1}\one.
\end{equation}
Select $w\notin Kv$ and take $B=vw^t-wv^t$, up to the nonzero
cofactor scalar. The range of $A_nB$ contains $d$. For a Jordan
link $E_{pq}$, the coefficient in \eqref{eq:cycleChord} is
proportional to
\begin{equation}\label{eq:linw}
 (a^tv)(w^tb)-(a^tw)(v^tb).
\end{equation}
If $a,b$ are independent, this is a nonzero linear functional of
$w$ precisely when $(a^tv,b^tv)\ne(0,0)$. The independence holds
for the nonforward edges used below: the interval is proper and
contains its two endpoints.

Here and below every Jordan block is represented with lower
subdiagonal links, and the coordinates may be permuted. For an
edge joining equal diagonal entries $\lambda$, formula
\eqref{eq:vdiag} yields
\begin{equation}\label{eq:intervalrelation}
 \frac12a^tv+b^tv=\lambda-\frac{\tau}{2n-1}.
\end{equation}
Indeed the left side assigns weight one to the two endpoints,
and weight one to every other coordinate; the sum is
$\sum_i v_i-(v_p+v_q)/2$. Substituting
$\sum_i v_i=\tau/(2n-1)$ and the common endpoint entries
$v_p=v_q=-\lambda+2\tau/(2n-1)$ gives the identity.

We consider all possible Jordan diagonals.
\smallskip

\noindent\emph{(i) $\tau\ne0$.}
List the Jordan blocks consecutively and use their ordinary lower
links $E_{i+1,i}$. Each relevant inclusive interval has just two
vertices. If $a^tv=b^tv=0$, the two equal endpoint entries of $v$
would sum to zero and so both vanish. The sum of the other entries
would vanish too, contradicting
$\sum_i v_i=\tau/(2n-1)\ne0$. For $n\ge3$ none of these lower
links is a forward cycle edge.

\smallskip
\noindent\emph{(ii) $\tau=0$ and $d\ne0$, with no nontrivial
zero-eigenvalue Jordan chain.}
Again use the consecutive Jordan ordering. Every required link
has nonzero eigenvalue $\lambda$; \eqref{eq:intervalrelation}
shows $a^tv$ and $b^tv$ cannot both vanish.

\smallskip
\noindent\emph{(iii) $\tau=0$, $d\ne0$, with at least one
nontrivial zero-eigenvalue Jordan chain.}
Two-color the vertices in every zero-eigenvalue Jordan chain as a
path, black and white, so each required zero link joins different
colors; assign any zero-eigenvalue singleton to either group.
Both colors occur. Select the leading endpoint of one nonzero-eigenvalue Jordan
chain as the single vertex in group $A$, with eigenvalue
$\mu\ne0$, and put the remaining nonzero-eigenvalue vertices
in group $B$, retaining their chain ordering. Since the total trace is zero, the sum of the
eigenvalues in $B$ is $-\mu$, so this group is nonempty. Arrange
the cyclic coordinates in the order
\begin{equation}\label{eq:blackwhite}
 \text{black zero vertices}\;|\; A\;|\;
 \text{white zero vertices}\;|\; B.
\end{equation}
A link with nonzero eigenvalue has nonzero functional by
\eqref{eq:intervalrelation}; its endpoints are not a forbidden
forward cycle edge. A zero Jordan link joins black to white,
and the inclusive cyclic interval from its source to its target
contains precisely one of $A,B$ among nonzero vertices. Thus
$\sum_{i\in I}d_i=\mu$ or $-\mu$. Now $v=-d$ and the zero
endpoints contribute nothing, so
$ a^tv=-2\sum_{i\in I}d_i=\mp2\mu\ne0$.
This deals simultaneously with several zero blocks and the
possible singleton vertices.

\smallskip
\noindent\emph{(iv) $d=0$.}
The target is nilpotent, possibly with several Jordan blocks.
Use the consecutive lower-Jordan ordering but choose $v=\one$
instead of \eqref{eq:vdiag}. The zero diagonal belongs to the
range of $A_nB$ automatically. For each required link the
inclusive interval has two vertices, so
$ a^tv=2$ and $b^tv=n-2$, and \eqref{eq:linw} is a
nonzero functional of $w$.

\smallskip
In each case we have finitely many nonzero linear functionals of
$w$ whose vanishing must be avoided, together with the proper
subspace $Kv$. Over the infinite field $K$, their union does not
exhaust $K^n$. Choose $w$ outside it, and select diagonal fixed
inputs realizing the decomposable bivector, as in
Lemma~\ref{lem:cycle}. Their coefficients on all required links
are nonzero. Solve \eqref{eq:cycleDiag} for an input $\diag(x)$
giving $d$ (with $x=0$ sufficient in (iv)), then add to \emph{that
same free input} the required $E_{pq}$ divided by their
coefficients in \eqref{eq:cycleChord}. The result is precisely the
permuted Jordan matrix $J$; the slice introduces no unwanted
matrix entries. Finally, conjugation gives the original target.
\end{proof}

\begin{corollary}\label{cor:minimal}
For every $n\ge3$, \eqref{eq:minimalclaim} holds.
\end{corollary}

\section{Intermediate matrix sizes: periodic path colors}
\label{sec:middle}
We now consider $r<n<2r$. Let $k=r-1$ and write
\begin{equation}\label{eq:kt}
 n=k+t+1,\qquad 1\le t\le k.
\end{equation}
First assume $k\ge3$. On the path with $n$ vertices, give
edge $e\to e+1$ the color $e\bmod k$. Put
\begin{equation}\label{eq:periodic}
\begin{split}
 F_j&=\sum_{\substack{0\le e<n-1\\e\equiv j\ (\mathrm{mod}\ k)}}
 E_{e,e+1},\\
 R_j&=\sum_{\substack{0\le e<n-1\\e\equiv j\ (\mathrm{mod}\ k)}}
 (-1)^{\lfloor e/k\rfloor}E_{e+1,e},\qquad 0\le j<k.
\end{split}
\end{equation}
Fix the inputs in the order
\begin{equation}\label{eq:periodorder}
 F_0,\ldots,F_{k-1},R_{k-1},\ldots,R_0
\end{equation}
and denote the final-variable slice by $L$.
Each term in the expansion chooses one forward and one reverse edge
of each of the $k$ colors. The colors $0,\ldots,t-1$ occur twice,
and the others once.

\subsection{Diagonal restriction and the connected-window rule}
A connected selection with the same forward and reverse edges must
be a window of $k$ consecutive edges. Every such window contains
each color once. The signs of the permutations of forward and
reverse colors cancel; the reversed-edge scalar signs have product
$(-1)^a$ when the window starts at $a$. By
Lemma~\ref{lem:pathloop},
\begin{equation}\label{eq:perdiag}
 L_{\rm diag}=\left(\sum_{a=0}^{t}
 \iota_a H\iota_a^t\right)\mathcal E_n.
\end{equation}
The Gram argument in Proposition~\ref{prop:positive} shows this
restriction is invertible: $H$ is positive definite over the reals,
and every coordinate belongs to one of the $t+1$ windows. Its
determinant is a nonzero integer.

We will use ordinary offsets as well as cyclic ones. For
$|h|<n$, put
\[
 \mathcal U_h=\Span\{E_{i,i+h}:0\le i,i+h<n\}.
\]
Because \eqref{eq:periodorder} contributes $k$ forward and $k$
backward unit steps, inserting $E_{u,u+h}$ gives output in
$\mathcal U_h$. Thus $L$ preserves \emph{ordinary} offsets,
and consequently respects their summands inside a cyclic offset.

A flow comparison across each path cut shows that, if the
coefficient from $E_{u,u+h}$ to $E_{i,i+h}$ is nonzero, the
multisets of colors on the oriented intervals beginning at $u$
and at $i$ of length $|h|$ must coincide. Removing complete cycles
of $k$ colors, a remaining nonzero proper color interval has a
unique start residue. Hence
\begin{equation}\label{eq:colorresidue}
 h\not\equiv0\pmod k\quad\Longrightarrow\quad i\equiv u\pmod k.
\end{equation}
This restriction is particularly strong since $n\le2k+1$:
for offsets two and three there are at most two positions in a
residue class.

For a coefficient along the diagonal of an ordinary-offset
restriction, input and output have the same endpoints, hence
forward and reverse selected edge sets agree. Encode the choice
of edge from duplicated color $j<t$ by $\epsilon_j=0$ for its
first occurrence and $\epsilon_j=1$ for its second. The edge
indicator along the path is
\begin{equation}\label{eq:binary}
 (1-\epsilon_0,\ldots,1-\epsilon_{t-1},
 \underbrace{1,\ldots,1}_{k-t},
 \epsilon_0,\ldots,\epsilon_{t-1}).
\end{equation}
When $t<k$, this pattern has one plus twice the number of
$0\to1$ rises of $\epsilon$ as its number of connected
components; the only permissible connected choices are the
ordinary windows. A single free chord cannot make one Euler
trail out of three or more components. When $t=k$, there is in
addition a two-component selection
$\epsilon=0^a1^{k-a}$ for each $1\le a<k$:
it takes one edge interval at the far left and another at the far
right, separated by a gap of $k$ edges. Such a selection cannot
be joined by a chord of length two if $k\ge3$, nor by a chord
of length three if $k\ge4$. We retain it for the long wrapping
offsets below.

\subsection{The nonwrapping offset-two restriction}
Consider $L|_{\mathcal U_2}$. By \eqref{eq:colorresidue},
its matrix splits into singleton blocks and $2\times2$ blocks
on the positions $i,i+k$ with $0\le i\le t-2$.
For a diagonal coefficient, the connected-window rule and
Lemma~\ref{lem:pathchord} show that all contributions have sign
$(-1)^i$, and their unsigned contributions are positive integers.
Some window contains both endpoints of $E_{i,i+2}$. For $k\ge3$
the multiplicity product in that window is at least two, so
\begin{equation}\label{eq:diagbound}
 |(L|_{\mathcal U_2})_{ii}|\ge2.
\end{equation}
For a pair $i,i+k$, the off-diagonal coefficient has absolute
value one. For example, to carry the free input
$E_{i+k,i+k+2}$ into output $E_{i,i+2}$, the cut equations force
the unmatched forward edges at $i,i+1$ and reverse edges at
$i+k,i+k+1$. Connectivity forces every edge in between. There
is then one Euler trail: forward from $i$ to $i+k$, take the
free chord, and return backward to $i+2$. Reversing this
construction gives the reverse coefficient, again of magnitude
one. The two-component exceptional selections cannot intervene
because their gap exceeds two.

Thus a paired restriction block has diagonal product of absolute
value at least four and off-diagonal product of absolute value
one; its determinant is nonzero. Every singleton is nonzero by
\eqref{eq:diagbound}. Hence
\begin{equation}\label{eq:U2}
 L(\mathcal U_2)=\mathcal U_2.
\end{equation}

\subsection{The wrapping offset-two restriction}
The remaining two units in $V_2$ have ordinary offset $2-n$.
It is convenient to compute the opposite offset $n-2$ and use
an anti-involution. Choose a diagonal $Z$ with
\begin{equation}\label{eq:Z}
 \frac{Z_{e+1,e+1}}{Z_{ee}}=(-1)^{\lfloor e/k\rfloor}
 \qquad(0\le e<n-1),
\end{equation}
and define $\theta(X)=ZX^tZ^{-1}$.
The map $\theta$ reverses products and interchanges the forward
and reverse colored inputs. Reversing the word of length $2k+1$
contributes $(-1)^k$, and reversing the $2k$ fixed inputs
contributes the same sign. Consequently
\begin{equation}\label{eq:theta}
 L\theta=\theta L,
\end{equation}
so opposite ordinary-offset restrictions are conjugate and have
the same invertibility status.

There are two positions on $\mathcal U_{n-2}$. Its restriction
is given by the following complete table (the displayed ordered
pairs are diagonal entries):
\begin{equation}\label{eq:wrapping}
\begin{array}{@{}lll@{}}
 \text{case}&\text{diagonal entries}&\text{off-diagonal entries}\\\hline
 2\le t<k,\ n\text{ even}&(4,-4)&0\\
 2\le t<k,\ n\text{ odd}&(-2,-2)&0\\
 t=1,\ k\text{ even}&(3,-3)&\text{each of magnitude}\le1\\
 t=1,\ k\text{ odd}&(-2,-2)&\text{each of magnitude}\le1\\
 t=k&(-(2k-1),-(2k-1))&0.
\end{array}
\end{equation}
We give the signed counts because the last row does not follow
from the ordinary connected-window rule alone.

For $2\le t<k$, color conservation prevents mixing between the
two positions. Applying \eqref{eq:pathchordtable} to the
connected windows gives respectively
$1+3(-1)^{n-2}$ and $-3+(-1)^{n-1}$.
For $t=1$, the two window contributions give
$\mathbf1_{2\mid k}+2(-1)^k$ and
$-2-\mathbf1_{2\mid k}$. A possible off-diagonal trail has a
turning point at one of $k$ positions and the signs alternate,
so its sum is $0$ or $\pm1$. In either case the determinant is
nonzero.

If $t=k$, the connected windows contribute $-2$ to each
of the two diagonal coefficients. For the extra selections
$\epsilon=0^a1^{k-a}$, the left and right bidirected components
have lengths $a$ and $k-a$ and the reverse-color sign is
$(-1)^{k-a}$. For the long chord $E_{0,2k-1}$, the local
starting sign at the left endpoint is positive, and the sign of
the right component at its penultimate vertex is
$(-1)^{k-a-1}$. Thus every two-end contribution has negative
sign after the reverse-color sign is included. Its magnitude
is two except for $a=k-1$, where it is one. For the other long
chord $E_{1,2k}$, the left local sign is negative and the right
endpoint sign is $(-1)^{k-a}$; again all contributions are
negative, with the unique magnitude-one case $a=1$. More explicitly, for $b=k-a$ the contributions of one
selection to the respective diagonal positions are
\begin{equation}\label{eq:uniform-wrap-selections}
 -(2-\mathbf1_{a=k-1}),\qquad
 -(2-\mathbf1_{a=1}),\quad 1\le a\le k-1.
\end{equation}
The reverse-color sign $(-1)^b$ is included in both expressions.
Summing the $k-1$ choices gives $-(2k-3)$ in each position.
Including the connected-window contributions $-2$ gives
$-(2k-1)$; off-diagonal entries vanish by
\eqref{eq:colorresidue}, since the two long-chord positions
have different residues modulo $k\ge3$. This proves the
last row of \eqref{eq:wrapping} uniformly in $k$, including
its signs.

Every restriction in \eqref{eq:wrapping} is invertible. By
\eqref{eq:theta}, so is the restriction on $\mathcal U_{2-n}$.
Combining it with \eqref{eq:U2} yields
\begin{equation}\label{eq:V2full}
 L(V_2)=V_2.
\end{equation}
For odd $n$, \eqref{eq:perdiag} and \eqref{eq:V2full}
now imply surjectivity by Lemma~\ref{lem:jordanpath}.

\subsection{The additional edge for even sizes}
Assume $k\ge4$. For the ordinary-offset-three restriction,
\eqref{eq:colorresidue} isolates the residue-zero block. Its
index set is $\{0\}$ if $t\le2$, and $\{0,k\}$ if $t\ge3$.
The corresponding matrix is
\begin{equation}\label{eq:grade3}
 \begin{cases}
 [-2],&t=1,\\{}[-4],&t=2,\\
 \begin{pmatrix}-5&\varepsilon\\\eta&5(-1)^k\end{pmatrix},
  &3\le t\le k,
 \end{cases}
 \qquad |\varepsilon|=|\eta|=1.
\end{equation}
For the row at $0$, the windows starting at $1,2,3$ contribute
$-2,-2,-1$ when present; the window at zero contributes zero,
since its endpoints have opposite parity. For the row at $k$,
the first three windows contribute $1,2,2$, with common sign
$(-1)^k$. The paired off-diagonal coefficients are the
unique Euler-trail values $\pm1$, as for ordinary offset two;
the two-end selections have too long a gap. This proves
\eqref{eq:grade3}. Every block is invertible: in the last
case its diagonal product has absolute value $25$, while its
off-diagonal product has absolute value $1$. Hence the slice
contains $E_{03}$, by choosing the appropriate free input
inside the isolated residue-zero block. For even $n$, the
Hamiltonian path \eqref{eq:evenpath} completes the argument.

\begin{proposition}\label{prop:middle}
For $r\ge5$ and $r<n<2r$, the image of $s_{2r-1}$ is
$\Mat_n(K)$. The same argument gives this conclusion for
$(r,n)=(4,5)$ and $(4,7)$.
\end{proposition}
\begin{proof}
For $k\ge4$ the diagonal and cyclic offset-two subspaces are
in the same fixed slice by \eqref{eq:perdiag} and
\eqref{eq:V2full}; for even $n$ it also contains $E_{03}$
by \eqref{eq:grade3}. Apply Lemma~\ref{lem:jordanpath}.
If $k=3$, the relevant intermediate odd sizes are $n=5,7$;
step two is Hamiltonian and the same diagonal and offset-two
arguments apply. The even intermediate size $n=6$ is treated
separately below.
\end{proof}

\section{Exceptional small matrix orders}\label{sec:seams}
The preceding uniform proofs leave $s_5$ on $\Mat_4,\Mat_5$
and $s_7$ on $\Mat_6$. We now give concrete slices or a block
construction for each. In the two fixed-slice cases, the complete
restriction matrices below are exact integer certificates; the
recurrences for regenerating every entry are given in
Appendix~\ref{app:verification}.

\subsection{\texorpdfstring{$s_5$ on $\Mat_5$}{s5 on M5}}
Use \eqref{eq:periodic} with $k=t=2$ and the fixed-input order
\eqref{eq:periodorder}. This case is not covered by the previous
proof, which assumed $k\ge3$. The restrictions to the diagonal
and cyclic step-two spaces, in the natural unit bases, are
\begin{equation}\label{eq:s5mat}
L_0=\begin{pmatrix}
2&-2&1&0&0\\
2&-8&4&-1&0\\
1&-4&10&-4&1\\
0&-1&4&-8&2\\
0&0&1&-2&2
\end{pmatrix},\qquad
L_2=\begin{pmatrix}
4&-1&1&0&0\\
-1&-4&-1&0&0\\
1&-1&4&0&0\\
0&0&0&-3&0\\
0&0&0&0&-3
\end{pmatrix}.
\end{equation}
Direct exact evaluation gives
\begin{equation}\label{eq:s5dets}
 \det L_0=900,\qquad\det L_2=-594.
\end{equation}
Both restrictions are invertible. Since step two is Hamiltonian
on five vertices, Lemma~\ref{lem:jordanpath} proves
\begin{equation}\label{eq:s5m5}
 \im s_5(\Mat_5(K))=\Mat_5(K).
\end{equation}

\subsection{\texorpdfstring{$s_7$ on $\Mat_6$}{s7 on M6}}
Take the Vandermonde core \eqref{eq:Vandermonde} with
$r=4$ and $q=6$. Its diagonal restriction has determinant
\begin{equation}\label{eq:s7d0}
 \det L_0=-31\,911\,936.
\end{equation}
The cyclic offset-two restriction is
\begin{equation}\label{eq:s7mat2}
 L_2=\begin{pmatrix}
 -16&-24&24&-4&0&0\\
 24&-132&144&-24&0&0\\
 -24&144&60&24&0&0\\
 4&-24&24&16&0&0\\
 0&0&0&0&-416&-36\\
 0&0&0&0&-36&-376
 \end{pmatrix}.
\end{equation}
The exact determinant of this restriction is
\[\det L_2=1\,838\,804\,889\,600.\]
For the parity-linking step the cyclic offset-three restriction is
\begin{equation}\label{eq:s7mat3}
 L_3=\begin{pmatrix}
164&44&0&0&0&0\\
44&432&100&0&0&0\\
0&100&52&0&0&0\\
0&0&0&164&-44&0\\
0&0&0&-44&432&-100\\
0&0&0&0&-100&52
\end{pmatrix},
\quad\det L_3=3\,776\,896\,843\,776.
\end{equation}
All three restrictions belong to the \emph{same} final-variable
slice, and all are invertible. Thus it contains every diagonal,
all of $V_2$, and in particular $E_{03}\in V_3$. The even-size
Hamiltonian path \eqref{eq:evenpath} gives
\begin{equation}\label{eq:s7m6}
 \im s_7(\Mat_6(K))=\Mat_6(K).
\end{equation}

\subsection{\texorpdfstring{$s_5$ on $\Mat_4$}{s5 on M4}: an explicit three-dimensional core}
We use an elementary $3\times3$ construction. Put
\begin{equation}\label{eq:s5core}
 A=E_{01},\quad B=E_{12},\quad C=E_{20},\quad
 D=\diag(-1,0,1)
\end{equation}
inside $\Mat_3(K)$. Direct expansion of the alternator yields
\begin{equation}\label{eq:s4endpoint}
 s_4(A,B,C,D)=\diag(-1,2,-1),
\end{equation}
and the slice image is the five-dimensional subspace
\begin{equation}\label{eq:W}
 W=\left\{
 \begin{pmatrix}\alpha&0&z\\u&\beta&0\\0&v&\alpha\end{pmatrix}
 :\alpha,\beta,u,v,z\in K\right\}.
\end{equation}
Indeed the specified matrix is obtained from the single free input
\begin{equation}\label{eq:Wpre}
 X=\diag\left(\frac{3\beta-4\alpha}{5},
 \frac{2\beta-\alpha}{5},0\right)
 +\frac u2E_{10}+\frac v2E_{21}-\frac z4E_{02}.
\end{equation}
Equations \eqref{eq:s4endpoint} and \eqref{eq:Wpre} follow by
direct evaluation.

\begin{lemma}\label{lem:Wclasses}
Every similarity class in $\Mat_3(K)$ meets $W$.
\end{lemma}
\begin{proof}
For a cyclic target with characteristic polynomial
$x^3-tx^2+qx-\delta$, choose $\alpha\in K$ solving
\[
 2t\alpha-3\alpha^2=q,
 \qquad \beta=t-2\alpha,
 \qquad u=v=1,\qquad z=\delta-\alpha^2\beta.
\]
The matrix in \eqref{eq:W} has the required characteristic
polynomial. Moreover
$\det[e_0,Me_0,M^2e_0]=1$, so it is cyclic and the
characteristic polynomial identifies its cyclic similarity class.
For noncyclic nonscalar classes the elementary representatives
$\diag(c,d,c)$ and $cI+E_{10}$, as appropriate to the Jordan
partition, lie in $W$. Scalars plainly do also.
\end{proof}

Embed \eqref{eq:s5core} into the upper-left $3\times3$ block of
$\Mat_4(K)$, with zero lower-right singleton. Because the
even-standard endpoint \eqref{eq:s4endpoint} is invertible, the
two off-diagonal rectangular maps are arbitrary, and the
upper-left block ranges over $W$. The lower-right output is
always zero. Thus this one slice has image
\begin{equation}\label{eq:3plus1}
 \left\{\begin{pmatrix}Y&b\\c&0\end{pmatrix}:
 Y\in W,\ b\in K^{3\times1},\ c\in K^{1\times3}\right\}.
\end{equation}

Let $T\in\Mat_4(K)$ be nonscalar. There exists $v$ such that
$v,Tv$ are independent; otherwise every vector is an eigenvector
and $T$ would be scalar. Choose $\phi\in(K^4)^*$ with
$\phi(v)=1$, $\phi(Tv)=0$. In a basis adapted to
$K^4=\ker\phi\oplus Kv$, the bottom-right coefficient of $T$
is zero. By Lemma~\ref{lem:Wclasses}, a similarity change on
$\ker\phi$ makes its $3\times3$ compression a member of $W$.
The other two blocks are unrestricted in \eqref{eq:3plus1}.
Hence every nonscalar $T$ is a single $s_5$ value, up to
conjugation.

To obtain scalars we need a short general observation.
\begin{lemma}[Traceful multilinear values]\label{lem:scalar}
Let $f$ be multilinear on $\Mat_n(K)$, with $K$ algebraically
closed of characteristic zero. If $f$ has an evaluation with
nonzero trace, then some evaluation of $f$ is a nonzero scalar
matrix. Consequently every scalar matrix is in its actual image.
\end{lemma}
\begin{proof}
Fix a primitive $n$th root of unity and the usual clock-and-shift
matrices $U,V$ with $UV=\zeta VU$. The $n^2$ matrices
$U^aV^b$, $a,b\bmod n$, form a basis of $\Mat_n(K)$ and are
homogeneous for the $\mathbb Z_n^2$ grading. Every product of
a fixed tuple of homogeneous basis elements lies in the same
graded component. By multilinearity, expanding a traceful
evaluation in this basis gives a homogeneous basis tuple on which
$f$ has nonzero trace. Every nonzero graded component has trace
zero, whereas the degree-zero component is $KI$. Thus this
particular value is a nonzero scalar. Rescaling one input gives
any scalar.
\end{proof}

The core \eqref{eq:s5core} is traceful: for example,
\eqref{eq:W} contains $\diag(1,0,1)$, which is a single
$s_5$ value on $\Mat_3(K)$. Embedding the same evaluation in
$\Mat_4(K)$ preserves its nonzero trace. Lemma~\ref{lem:scalar}
therefore gives all scalar targets, completing
\begin{equation}\label{eq:s5m4}
 \im s_5(\Mat_4(K))=\Mat_4(K).
\end{equation}

\section{Proof of the odd-standard assertion}\label{sec:assembly}
\begin{proof}[Proof of the odd assertion of Theorem~\ref{thm:main}]
If $n<r$, then $2r-1\ge2n+1$. By Amitsur--Levitzki,
$s_{2n}$ is an identity on $\Mat_n(K)$; the standard recursion
expresses every higher alternator through lower ones, proving
$s_{2r-1}=0$.

For $r=1$, the claim is $s_1(X)=X$. For $r=2$ and $n=1$,
the alternating cubic vanishes. For $n\ge2$, the published
cubic image-linearity result \cite[Theorem~3.5]{Vitas2026}
implies that $\im s_3$ is a vector subspace. The substitution
$s_3(I,A,B)=[A,B]$ gives every traceless matrix by the classical
commutator theorem \cite{AlbertMuckenhoupt1957}, and a traceful
$2\times2$ evaluation,
embedded into $\Mat_n(K)$, gives a value outside that subspace.
For example, $s_3(E_{00},E_{01},E_{10})=\diag(2,1)$
on $\Mat_2(K)$. Since the traceless subspace has codimension one, cubic
image-linearity now gives the whole $\Mat_n(K)$.

Suppose $r\ge3$. The case $n=r$ follows from
Corollary~\ref{cor:minimal}. The range $n\ge2r$ follows from
Corollary~\ref{cor:large}. For $r<n<2r$, when $r\ge5$ apply
Proposition~\ref{prop:middle}; when $r=4$, it covers $n=5,7$
and \eqref{eq:s7m6} covers $n=6$. When $r=3$, the sizes
$n=4,5$ are \eqref{eq:s5m4} and \eqref{eq:s5m5}.
The cases are partitioned as follows:
\begin{center}
\small
\begin{tabular}{@{}ll@{}}
\toprule
Degree parameter & Surjective sizes and argument \\
\midrule
$r=1$ & all $n\ge1$; identity map\\
$r=2$ & all $n\ge2$; cubic theorem and direct values\\
$r=3$ & $n=3$ cycle, $n=4,5$ cases, $n\ge6$ transfer\\
$r=4$ & $n=4$ cycle; $n=5,7$ periodic; $n=6$ case; $n\ge8$ transfer\\
$r\ge5$ & $n=r$ cycle; $r<n<2r$ periodic; $n\ge2r$ transfer\\
\bottomrule
\end{tabular}
\end{center}
Together with the identity branch $n<r$, this exhausts every
pair of positive integers $(r,n)$.
\end{proof}

\begin{remark}
The large and periodic ranges use a fixed linear slice, whereas the
minimal-size construction chooses a slice for the prescribed Jordan target.
\end{remark}

